\section{Experiments}

\subsection{Research Questions}

The experiments are organized around four questions:
\begin{itemize}
  \item \textbf{RQ1: Do LLM-generated evaluation criteria exhibit systematic
  blind spots?} This is tested against human-gold evaluation dimensions with BSC and
  confidence intervals.
  \item \textbf{RQ2: When do the evidence gates warrant a dimension-level
  coverage-change or recovery statement without merely increasing length?} This
  requires reporting coverage, blind-spot rate, redundancy, hallucination,
  dimension-transition audits, and criteria-budget curves.
  \item \textbf{RQ3: Is a BSC coverage change supported by downstream judging utility?}
  Generated criteria are used to score chosen-vs-rejected or multi-candidate
  benchmark examples.
  \item \textbf{RQ4: Which reward components matter?} We ablate redundancy,
  validity/hallucination filtering, verifier filtering, and the RL stage.
\end{itemize}

\subsection{Datasets and Isolation}

The main BSC benchmark is RubricBench \texttt{test\_main}, a 497-example
human-gold holdout with 491 unique queries that is not used for SFT, proxy criteria elicitation, or reward tuning.
It is also not used for prompt selection, hyperparameter tuning, checkpoint
selection, or verifier calibration; it is opened only for the final hard-gold
BSC evaluation and its audits.
RubricBench \texttt{train\_seed} and ResearchRubrics \texttt{train\_seed}
provide 560 human-gold seed examples; these seed examples are used only for
small-scale initialization and protocol calibration, not as the main holdout.
Proxy-gold training data is used only for scaling supervision and comes from
non-overlapping query pools: RewardBench \texttt{sft\_proxy\_train} (1,785
after filtering against RubricBench \texttt{test\_main} and downstream
RewardBench, JudgeBench, and RewardBench-2 holdouts),
IFBench (300), WritingBench (1,000), HealthBench (5,000), and BEIR/NQ (3,452),
after multi-teacher criteria elicitation and verifier filtering. Downstream utility is
evaluated on non-overlap held-out RewardBench, RewardBench-2, and JudgeBench
splits. The corresponding audit files contain 597 RewardBench records
(551 normalized unique queries), 1,815 RewardBench-2 records (1,814 raw stripped
unique queries; 1,813 normalized unique queries), and 620 JudgeBench records
(528 normalized unique queries). This separation lets us write the paper around hard-gold evaluation
while still using proxy-gold data for scale.
Any paper-facing method-result claim must therefore bind the exact training
inputs, proxy-train filtering reports, zero-overlap holdout audits, and
evaluation outputs. The binding is content-addressed: filtering reports expose
holdout/input/output SHA256 hashes, final contamination audits expose holdout
and training-artifact hashes, and the evidence matrix checks these hashes
across the filtering chain and against the verifier-filtering,
SFT/proxy-gold, and GRPO construction reports. The verifier-filtering report
must bind the raw teacher criteria, filtered teacher criteria, Meta-Verifier
provider, SFT-data preflight report, and meta-verifier budget report by SHA256;
its filtered-output digest must match the proxy-gold construction input digest.
An isolated table row is not considered evidence without the corresponding
non-overlap audit and hash-consistent provenance. A report with
\texttt{overlap\_query\_count == 0} but
\texttt{artifact\_status=blocked} is treated as not auditable rather than as
clean isolation, because it has not inspected the exact SFT, proxy-gold, and
GRPO artifacts used by the trained row.

\subsection{Minimal Motivation Experiment}

The first experiment measures whether a single model misses human evaluation
dimensions on hard-gold criterion data. The required evidence is the BSC summary,
per-item diagnostics, threshold sweep, and claim-evidence matrix generated by
the minimal claim pipeline. This minimal motivation gate is safe for Section 2
only; it is not the full real-run readiness matrix and does not unlock any
trained-method row. Under that limited gate, on 100 RubricBench hard-gold
examples, the base evaluation-criteria policy obtains 0.3692 coverage and
0.6308 blind-spot rate, with a 95\% CI of 0.5785--0.6823 for the blind-spot
rate.

\subsection{Main Matrix}

The main comparison evaluates the base evaluation-criteria policy, strong API
teachers, the SFT-only policy, and the full SFT+RL policy. The intended metrics
are coverage, blind-spot rate, redundancy, hallucination, and downstream
chosen-vs-rejected accuracy. The main table should not be interpreted as
complete until SFT-only, SFT+GRPO, and downstream rows pass the claim-evidence
gates.

For the central method claim, these metrics are conjunctive rather than
independent: a paper-facing method-result claim requires a reportable hard-gold
semantic coverage change under the fixed protocol, no material degradation in redundancy or
hallucination, and held-out downstream judge-utility support. A row with a
numerically stronger BSC value alone is treatable only as metric evidence, not
as judge-utility support.

\subsection{Claim Ladder and Downgrade Rules}

The paper uses a claim ladder rather than a single all-or-nothing success
criterion. The frozen 100-example hard-gold diagnostic is a \emph{motivation}
claim: it supports the existence of systematic evaluation blind spots for a
base evaluation-criteria policy, but it does not by itself support a trained
method claim. A hard-gold BSC coverage change with stable redundancy and
hallucination is a \emph{metric-support} claim: it reports how many
human-gold dimensions are missed under the fixed semantic protocol. It becomes \emph{method-support}
only when the SFT-only versus SFT+GRPO comparison and reward-component
ablations support attributing the observed coverage change to the RLVR reward
rather than to length or duplicate criteria. It becomes \emph{judge-utility support} only
when the held-out RewardBench, RewardBench-2, and JudgeBench rows are
paper-eligible under the API/model scorer and SHA-bound budget contract.

If any higher gate fails, the manuscript downgrades the corresponding sentence
instead of carrying the stronger conclusion forward. Without downstream
support, the result is reported as metric-only BSC evidence. Without C12
dimension-transition support, the result is reported as an aggregate coverage
change rather than dimension-level recovery. Without C14 SFT-only versus SFT+GRPO
support, the result is reported as proxy-gold supervision evidence rather than
RLVR evidence. Without clean C0 contamination and provenance gates, no trained
method row is paper-facing regardless of its metric values.

\begin{table*}[t]
\centering
\begin{tabular}{p{0.12\linewidth}p{0.30\linewidth}p{0.26\linewidth}p{0.22\linewidth}}
\toprule
Gate & Required evidence & Claim unlocked & If missing \\
\midrule
C0 & Zero-overlap hard-gold, proxy-train, and downstream audits with
SHA-bound provenance & Any trained-method row can be paper-facing & Report only
non-training diagnostics \\
C2--C3 & Hard-gold BSC, redundancy, hallucination, confidence intervals,
and threshold robustness under one fixed protocol & Metric-support for
coverage changes & Do not write aggregate method conclusions \\
C4/C9/C10 & RewardBench, JudgeBench, and RewardBench-2 utility under the
fixed scorer and budget contract & Judge-utility support & Keep BSC results
metric-only \\
C7 & Reward-component and verifier ablations & Evidence supporting attribution
to reward components rather than verbosity & Treat as an uncontrolled
training comparison \\
C12 & Query-aligned per-gold dimension transition audit & Dimension-level
recovery evidence & Write aggregate coverage change only \\
C13 & Point-level semantic-space CSV, JSON summary, and rendered figure with
cluster-coverage checks & Mechanism visualization & Treat the figure as
illustrative only \\
C14 & SFT-only versus SFT+GRPO under the same protocol & RLVR-stage support &
Report proxy-gold supervision only \\
\bottomrule
\end{tabular}
\caption{Evidence gates used to keep empirical conclusions aligned with
available artifacts. The gates are conjunctive for the central method claim:
coverage, redundancy, hallucination, downstream utility, ablations,
dimension-level transitions, and visualization must all be supported before the
manuscript states a broad trained-method conclusion.}
\label{tab:evidence_gates}
\end{table*}

\IfFileExists{tables/main_table.tex}{\input{tables/main_table}}{
\typeout{Main results table not included because tables/main_table.tex is missing.}
}

\subsection{Ablations}

The minimum AAAI-facing ablation set is:
\begin{itemize}
  \item \textbf{No redundancy penalty}: tests whether BSC coverage changes come from
  repetitive paraphrases rather than new dimensions.
  \item \textbf{No hallucination/validity term}: tests whether coverage reward
  encourages unverifiable or query-ungrounded criteria.
  \item \textbf{No verifier filtering}: tests whether proxy-gold quality
  control matters for SFT and RL stability.
  \item \textbf{SFT-only vs. SFT+GRPO}: tests whether RLVR changes
  dimension-level coverage beyond supervised imitation of teacher criteria.
  \item \textbf{Single teacher vs. multi-teacher union}: tests whether teacher
  diversity is supported as a proxy-gold construction choice before RL.
\end{itemize}
The relevant RL-stage claim test is not whether RL produces longer criteria
lists. The RL-stage claim is safe only if the hard-gold evidence shows a
reportable semantic coverage change while redundancy and invalidity remain within the configured gates and
the mean generated-criteria count, generated-to-gold ratio, and coverage per
generated criterion remain bounded under the length-control audit. The
SFT-only and SFT+GRPO rows must also share the same BGE embedding model,
coverage and redundancy thresholds, verifier-backed hallucination source, and
bootstrap-CI protocol; otherwise we treat the stage comparison as not
reportable.

Each ablation has a specific failure mode. If removing the redundancy penalty
raises coverage only by increasing duplicate criteria, the change is treated as
unsupported for dimension-level recovery. If removing the validity term or
verifier filtering raises coverage while hallucination rises, the method is
memorizing or inventing criteria rather than producing faithful dimensions. If
SFT-only matches SFT+GRPO under the same protocol, the RLVR stage should be
reported as unsupported for the observed coverage change. If the multi-teacher union does
not beat the strongest single teacher, proxy-gold scaling should be described
as an uncontrolled data-construction choice rather than an empirically supported data-construction
advantage.
Offline reward re-scoring on the same model outputs is used only as a diagnostic
for the reward formula; the paper-facing no-redundancy, no-validity,
no-verifier, and coverage-only rows require separately trained variants
evaluated on the same hard-gold protocol.
The trained-ablation completion artifact must enumerate
\texttt{full}, \texttt{no\_red}, \texttt{no\_valid},
\texttt{no\_verifier}, and \texttt{cov\_only}; it keeps the verifier enabled
for the full and no-redundancy variants, disables both the validity weight and
verifier reward for the no-validity variant, keeps validity and redundancy
weights while disabling only verifier reward for the no-verifier variant, and
disables the validity weight, redundancy weight, and verifier reward for the
coverage-only variant. The no-verifier filtering ablation is reportable only if
verifier filtering does not increase hallucination and does not collapse
hard-gold semantic coverage relative to the raw proxy-criteria variant.

\IfFileExists{tables/rl_stage_ablation_table.tex}{
\input{tables/rl_stage_ablation_table}
}{}

\subsection{Blind-Spot Map and Budget Analysis}

To test whether blind spots are systematic rather than random misses, we report
category-level blind-spot rates over the fixed Blind-Spot Attribution taxonomy.
The resulting map is intended to show which categories of human evaluation
dimensions remain uncovered for each model. We also evaluate a criteria budget
curve by truncating
each generated criteria set to $K \in \{3,5,8,10,15\}$ items and recomputing BSC.
This controls for the possibility that a policy raises coverage only by
eliciting more criteria. For trained policies, we additionally report
dimension-level transition rates: the fraction of baseline blind spots that
become covered by the trained policy, and the fraction of previously covered
dimensions that are lost. These metrics are computed at the gold-dimension
level rather than from aggregate averages: a human-gold criterion counts as
recovered only if it is missed by the baseline and matched by the candidate
under the same semantic threshold. The corresponding audit artifact is produced by
\texttt{transition\_summary.json}, \texttt{transition\_per\_item.csv},
\texttt{transition\_by\_category.csv}, and
\texttt{transition\_gold\_items.jsonl}, which also report lost dimensions that
were covered by the baseline but missed by the candidate. The
dimension-recovery claim is reportable only when the recovered/lost dimension
audit shows recovered human-gold dimensions exceeding lost dimensions, baseline
and candidate rows align by query with no unmatched records, use the same BGE embedding model and
$\tau=0.75$ coverage threshold, and respect verifier \texttt{valid\_flags}.

\subsection{Downstream Judge Utility}

BSC alone is not sufficient for a strong paper claim: reviewers may argue that
it optimizes a self-defined metric. We therefore evaluate whether generated
criteria help an evaluator choose the preferred response on downstream
benchmarks. For pairwise data, each candidate answer is scored against the
generated criteria set and the higher-scoring answer is selected. For
multi-candidate data, the criteria-scored top candidate is compared against the
gold label. The paper-facing scorer should be a fixed API/model scorer with
structured JSON outputs; deterministic keyword scoring is retained only for
smoke tests. The downstream table must include accuracy, tie rate, and margin
on RewardBench holdout, RewardBench-2, and JudgeBench.
Downstream utility is therefore the external validity check for BSC: if BSC
is numerically stronger but downstream accuracy, tie rate, or margin is not
supported under the API scorer, the manuscript should report the result as
metric-only evidence and not as judge-utility evidence.
Each downstream summary also records the exact joined evaluation input,
current input/provider SHA256 hashes, benchmark format, scorer provider,
budget report, scorer-call contract (pairwise versus multi-candidate), and a
\texttt{paper\_claim\_eligible} flag. The evidence matrix compares these
summary hashes against the budget-report contract hashes, so matching paths
alone cannot make stale API-budget reports or keyword smoke runs paper-facing
utility evidence.

\IfFileExists{tables/downstream_utility_table.tex}{
\input{tables/downstream_utility_table}
}{}

\subsection{Semantic-Space Visualization}

To audit the proposed semantic-coverage mechanism, we embed generated and
human-gold evaluation dimensions and project them into a two-dimensional
semantic space. The
paper-facing matrix requests UMAP projection and records the actual projection
used in the JSON summary, falling back to deterministic PCA only when optional
visualization dependencies are unavailable. The asset builder outputs
point-level CSV, JSON summary, and SVG/PDF figures with human-gold criteria
shown as hollow points and generated criteria as filled points, colored by
blind-spot category. The visualization audit tests whether SFT+GRPO criteria
show a reportable nearest-gold cluster-coverage change relative to SFT-only
criteria, where gold clusters are deterministic connected components of
human-gold dimensions under the same embedding protocol and coverage threshold
used by BSC. It also tests whether the redundancy ablation is consistent with
local-collapse evidence. This visualization is not a substitute for BSC, but it
makes the proposed mechanism auditable only under C13: a mechanism claim is
reportable only if the point-level evidence supports nearest-gold coverage
change without local collapse, rather than merely showing more near-duplicate
generated criteria. The evidence matrix therefore gates this figure on method-level gold-category
coverage, nearest-gold category coverage, nearest-gold gold-cluster coverage,
nearest-gold cluster-distribution entropy, generated-criteria dispersion, and
nearest-gold similarity rather than only on the presence of rendered image
files. If the figure looks visually broader but the point-level CSV and JSON
summary do not show nearest-gold cluster coverage change, non-collapsed cluster
distribution, and non-collapsed dispersion, we treat the visualization as
illustrative only, not as mechanism evidence.
The reproducible asset command is
\texttt{scripts/build\_semantic\_space\_visualization.py}.

\IfFileExists{figures/semantic_space.pdf}{
\begin{figure}[t]
  \centering
  \includegraphics[width=\linewidth]{figures/semantic_space.pdf}
  \caption{Semantic-space projection of human-gold and generated evaluation
  dimensions. Human-gold dimensions are plotted as hollow markers and generated
  dimensions as filled markers; the C13 audit tests whether SFT+GRPO has a
  reportable nearest-gold region-coverage change relative to SFT-only under
  the same embedding protocol, subject to point-level CSV/JSON evidence.}
  \label{fig:semantic_space}
\end{figure}
}{}

\IfFileExists{tables/ablation_table.tex}{\input{tables/ablation_table}}{}

\IfFileExists{tables/teacher_union_ablation_table.tex}{
\input{tables/teacher_union_ablation_table}
}{}

\IfFileExists{tables/verifier_filter_ablation_table.tex}{
\input{tables/verifier_filter_ablation_table}
}{}

\IfFileExists{tables/dimension_transition_table.tex}{
\input{tables/dimension_transition_table}
}{}
